\documentclass[12pt]{article}
\usepackage[utf8]{inputenc}
\usepackage[T1]{fontenc}
\usepackage{hyperref}
\title{Koro Nulu Language Resources for K-12 Curriculum Development}
\author{Compiled by Claude Sonnet 5 (claude-sonnet-5) for Andr\'as Kornai}
\date{}
\begin{document}
\maketitle
\section{Introduction}
Koro Nulu (also Koro Ija, Koro Afiki; ISO 639-3 vkn) is an unclassified Plateau language of Tafa LGA, Niger State, Nigeria. Estimates are about 4,000 to 4,700 speakers in Ija Koro, Kabwei, Dogon Kurmi and Shere. Glottolog rates it vigorous. Youth may be shifting toward Hausa or English. Very little published material exists.
\section{Primary Educational Resources}
None found. No primers, textbooks or school programmes were found, and Joshua Project reports no written form.
\section{Linguistic and Reference Materials}
\begin{itemize}
\item Wikipedia entry: \url{https://en.wikipedia.org/wiki/Koro_Nulu_language}
\item Glottolog entry: \url{https://glottolog.org/resource/languoid/iso/vkn}
\item A 2021 sociolinguistic profile of Koro Nulu and Koro Zuba in the Journal of Language Survey Reports was mentioned in search results, but no direct URL was seen, so it is not listed as a resource.
\end{itemize}
\section{Religious and Community Resources}
Joshua Project reports that the community is almost entirely Muslim and that Bible translation has begun: \url{https://joshuaproject.net/people_groups/21981}
\section{Digital and Audio Resources}
Global Recordings Network has a page for the language; it lists recordings for the related variety Koro: Adunu but none confirmed for Koro Nulu itself: \url{https://globalrecordings.net/en/language/vkn}
\section{Curriculum Development Context}
Speakers are reported to be positive about the language and want it developed. Orthography, literacy materials and community partnerships would have to be created first.
\end{document}